% TOP_PROBLEM: {"id": 2536, "title": "Kourovka Notebook Problem 21.27", "category_id": 17, "category": {"id": 17, "name": "group_theory", "display_name": "Group Theory", "description": "Problems about groups, group actions, representations, and related algebraic structures.", "slug": "group-theory", "order_index": 17, "created_at": "2026-07-31T15:26:25.671Z"}, "status": "open", "problem_number": "KOU-21.27", "set_id": 10, "statusQualification": "The nonabelian-simple convention of the derangement conjecture is explicit; literal inclusion of the regular cyclic group of order two gives DD={1}."}
% FIRST_POSTED: 2026-10-01
% ENTRY_KIND: statement_only
% UnsolvedMath ID 2536; source catalogue code KOU-21.27
% !TeX program = lualatex
% Definition and sources only; this file is not an LLM solution attempt.
\documentclass[11pt,a4paper]{article}
\usepackage[margin=25mm]{geometry}
\usepackage{fontspec}
\setmainfont{lmroman10-regular.otf}[BoldFont=lmroman10-bold.otf,
 ItalicFont=lmroman10-italic.otf,BoldItalicFont=lmroman10-bolditalic.otf]
\usepackage{amsmath,amssymb,mathtools}
\usepackage{unicode-math}
\setmathfont{latinmodern-math.otf}
\AtBeginDocument{\let\setminus\smallsetminus}
\usepackage{xurl,hyperref}
\hypersetup{colorlinks=true,urlcolor=blue}
\setcounter{secnumdepth}{0}
\setlength{\parindent}{0pt}
\setlength{\parskip}{5pt}
\setlength{\emergencystretch}{3em}
\begin{document}
\section*{Kourovka Notebook Problem 21.27}\label{2536}
{\small UnsolvedMath ID 2536 · KOU-21.27 · Group Theory · Kourovka Notebook - New Problems, Issue 21}
\subsection{Definitions and mathematical statement}
Let $G\le\operatorname{Sym}(\Omega)$ be a finite nonabelian simple group
acting transitively on the finite set $\Omega$. Being a permutation
subgroup specifies a faithful action: only the identity fixes every
point. Transitivity means that for all $\alpha,\beta\in\Omega$ some
$g\in G$ satisfies $g\alpha=\beta$. Simplicity means that $G$ has no
normal subgroup other than $1$ and itself.

An element $g$ is a derangement in this particular action if it fixes
no point of $\Omega$. Write
\[
 D=\{g\in G:g\alpha\ne\alpha\text{ for every }\alpha\in\Omega\}.
\]
Does the set product $DD=\{ab:a,b\in D\}$ equal $G$ for every such
action? In explicit terms, for every $g\in G$ must there be derangements
$a,b$ with $g=ab$?

The two factors are not required to be distinct and may depend on $g$.
The identity is included among the target elements; its representation
would use inverse derangements. The property is action-dependent, since
the same abstract simple group can have different transitive permutation
representations and different sets $D$.

Every faithful transitive action of every finite nonabelian simple group is included;
primitivity of the action or a restriction on point stabilizers is not
assumed. The claim concerns products of exactly two derangements, rather
than merely generation of $G$ by derangements or representation by an
unbounded number of factors.

The nonabelian convention is explicit here. If cyclic simple groups are
included literally, the regular action of $C_2=\{1,a\}$ has
$D=\{a\}$ and $DD=\{1\}$, an elementary exception to the printed wording.
\subsection{Short English statement}
In every finite nonabelian simple transitive permutation group, is each element a product of two fixed-point-free permutations from that group?
\subsection{Sources}
\begin{itemize}
\item[S1] ULAM.AI and the UnsolvedMath Contributors, supplied problems.json, record 2536 (KOU-21.27), Kourovka Notebook - New Problems, Issue 21. Catalogue identity and source transcription; CC BY 4.0.
\url{https://huggingface.co/datasets/ulamai/UnsolvedMath}.
\item[S2] The Kourovka Notebook, 21st edition, new problems, Problem 21.27. Expanded formulation of the supplied catalogue statement.
\url{https://arxiv.org/pdf/1401.0300v47}.
\item[S3] M. Larsen, A. Shalev and P. H. Tiep, Products of derangements in simple permutation groups, Forum of Mathematics, Sigma10 (2022). Primary literature for the simple-group derangement problem.
\url{https://doi.org/10.1017/fms.2022.69}.
\end{itemize}
\end{document}
